\section{Background}

The emergence of large language model (LLM)-driven multi-agent systems has introduced a fundamental architectural challenge: managing the lifecycle, lineage, and accountability of agents that spawn, delegate to, and depend on one another across distributed execution environments. As autonomous agent systems scale from single-task workflows to open-ended, recursively decomposing hierarchies, the need for principled governance of agent relationships becomes not merely operational but scientific. This section situates the Recursive Spawning model within five interconnected research threads: agent lifecycle management, accountability and provenance, delegation chain semantics, hierarchical task decomposition, and the BX3 Framework's approach to immutable parent chains.

\subsection{Agent Lifecycle Management in Multi-Agent Systems}

The management of agent lifecycles — from creation and initialization through execution, delegation, and termination — has become a central concern in both academic and industrial AI research. Traditional software lifecycle management provides a useful analogue: agents, like processes, require structured spawning, resource allocation, state monitoring, and graceful teardown. However, unlike operating system processes, AI agents are non-deterministic, stateful, and capable of modifying their own execution context, which complicates lifecycle governance substantially.

Recent work has formalized the task lifecycle in multi-agent systems. In \cite{agentlifecycle}, the authors derive 16 host-agent properties and 14 task-lifecycle properties expressed in temporal logic, covering states from task creation through completion, error handling, and delegation. This formalization reveals that most contemporary frameworks conflate execution with orchestration, providing minimal structure for reasoning about agent relationships over time. Similarly, \cite{agentrm} identifies two critical failure modes in multi-agent LLM deployments: scheduling failures (zombie agents, rate-limit cascades) and context degradation leading to agent amnesia. Their OS-inspired resource manager, AgentRM, demonstrates that treating agent slots as system resources enables substantial improvements in throughput and reliability, but stops short of addressing the provenance or accountability implications of dynamic agent creation and delegation.

The ALTK framework \cite{altk} provides modular middleware components for lifecycle intervention — post-request conditioning, pre-tool validation, post-output checking — but does so within a single-agent paradigm. The extension to multi-agent scenarios, where one agent's lifecycle events cascade into another's operational context, remains under-addressed. Protocols such as MCP \cite{mcp} and A2A \cite{a2a} standardize tool access and agent-to-agent communication, yet provide no native mechanism for tracking which agent spawned which, under what authority, and with what residual accountability. AGENTHIVE \cite{agenthive} addresses dynamic, delegation-first topologies where any agent can spawn sub-agents recursively, forming task-adaptive structures (trees, forests, star topologies) without pre-defined graphs. This delegation-as-first-class-primitive model directly motivates the Recursive Spawning architecture, where parent-child relationships are structural, not incidental.

\subsection{Accountability and Provenance in Distributed AI Systems}

Accountability in distributed AI systems requires the ability to trace decisions, actions, and outputs back to the agents and human principals responsible for them. Unlike traditional distributed systems where causality is deterministic and logs are append-only, LLM-driven agents introduce probabilistic outputs, context-dependent reasoning, and cross-tool epistemic chains where the provenance of a decision may span multiple agents, external tools, and model inference calls.

The shift from data lineage to provenance governance in AI systems is examined in \cite{forbesprov}, which argues that generative AI outputs are built from learned patterns across opaque training corpora, making exact source attribution fundamentally weaker than in rule-based systems. This insight has motivated new frameworks for AI-specific provenance. PROV-AGENT \cite{provegent} extends W3C PROV with concepts from MCP and data observability to capture AI-agent interactions — prompts, responses, decisions — and link them to the broader workflow and downstream outcomes. Its near real-time provenance capture system enables transparency, traceability, and reproducibility across heterogeneous environments.

More formal approaches have emerged from the IETF. The Verifiable AI Provenance Framework (VAP) \cite{vap} defines an architectural framework for evidentiary-grade provenance using existing security primitives (SCITT, RATS, COSE) to produce tamper-evident, cryptographically verifiable decision trails. Warrant Certificate Authorities (WCA) \cite{wca} extend this further by certifying data sources rather than data content, preventing what the authors term semantic laundering across tool boundaries. WCA introduces provenance layers with reference-monitor properties (complete mediation, tamper-resistance, verifiability) analogous to PKI trust models. The TIBET protocol \cite{tibet} specifies four-dimensional provenance tokens — capturing content (ERIN), dependencies (ERAAN), environmental context (EROMHEEN), and intent (ERACHTER) — linked via parent identifiers to form complete audit trails.

These provenance frameworks share a common deficiency: they address the evidence of decisions after the fact, but do not constrain the structural relationships between agents that produce those decisions. A parent chain that can be retroactively modified or pruned severs the evidentiary chain regardless of how detailed the individual tokens are. The Recursive Spawning model addresses this at the architectural level by making parent relationships immutable from the moment of spawning.

\subsection{Fixed vs. Mutable Delegation Chains}

Delegation — the act of one agent transferring a task or sub-task to another — is the foundational operation of multi-agent orchestration. The semantics of delegation, however, vary dramatically across implementations. In fixed delegation systems \cite{fixeddelegation}, an agent may spawn sub-agents with constrained authority, but the parent-child relationship is established once and persists without modification. Mutable delegation systems \cite{mutabledelegation} permit the re-parenting, pruning, or redirection of agents after creation, which introduces significant accountability hazards: an agent's lineage can be altered retroactively, obscuring the principal who originally authorized its actions.

The theoretical distinction matters for governance. Fixed delegation chains preserve \textit{accountability continuity} — the ability to trace any agent's decisions back through its ancestry without interruption. Mutable chains introduce \textit{accountability gaps} where a sub-agent may lose reference to its originating principal, either through re-parenting or through orphaned execution where the parent has been terminated or pruned.

StackPlanner \cite{stackplanner} demonstrates a centralized hierarchical multi-agent system where high-level coordination is separated from subtask execution, with explicit task-level memory management to prevent context bloat. However, its architecture is manager-centric: sub-agents are created and owned by a central planner, not by other agents in the hierarchy. This design simplifies accountability at the cost of flexibility. AGENTHIVE \cite{agenthive} inverts this by making delegation a first-class primitive available to any agent, enabling decentralized control without a fixed central orchestrator. The Recursive Spawning model extends this further by imposing immutability constraints on the parent chain, preserving the flexibility of decentralized delegation while maintaining the accountability continuity that mutable chains sacrifice.

The governance implications of mutable delegation are examined in GAAT \cite{gaat}, which provides real-time, closed-loop enforcement in multi-agent systems. Their Governance Telemetry Schema extends OpenTelemetry with governance attributes, and their enforcement bus achieves sub-200 ms detection latency with 98.3\% violation prevention rates. Yet GAAT operates at the telemetry layer — it observes and enforces but does not structurally prevent delegation chains from being altered after the fact. Immutable parent chains provide the structural precondition that makes telemetry-driven governance effective, because the chain of accountability is preserved rather than being subject to retroactive mutation.

\subsection{Hierarchical Task Decomposition in AI}

Hierarchical task decomposition — the process of recursively breaking down high-level goals into executable sub-tasks — is foundational to planning and reasoning in AI agents. Classical planning paradigms, particularly Hierarchical Task Networks (HTN) \cite{htn}, formalize this with task networks where compound tasks decompose into primitive actions via methods. Modern LLM-based agents extend this with dynamic, on-the-fly decomposition that adapts to the executor's capabilities and the environmental context.

ADaPT \cite{adapt} introduces as-needed decomposition where task splitting occurs only when required, recursively adapting to both task complexity and the executor LLM's capabilities. ReAcTree \cite{reactree} builds hierarchical agent trees with explicit control flow nodes, demonstrating that dynamic tree construction with episodic and working memory enables robust long-horizon planning — achieving nearly double the success rate of flat ReAct baselines on long-horizon embodied tasks. HiPlan \cite{hiplan} combines global milestone-level planning with local step-wise hints, demonstrating that hierarchical guidance improves both robustness and adaptability in changing environments.

STRUCTUREDAGENT \cite{structuredagent} introduces AND/OR tree search for hierarchical web agents, showing that explicit planning structures with dynamic tree search improve performance on complex multi-page tasks where greedy termination and limited in-context memory are primary failure modes. Task-Decoupled Planning (TDP) \cite{tdp} further decouples planning from execution via a DAG-based sub-goal structure, reducing token costs by up to 82\% while improving error isolation. HiMAC \cite{himac} provides a macro-micro learning framework that separates planning from execution, demonstrating that explicit hierarchy — not merely model scale — is the key to robust long-horizon decision-making.

These decomposition frameworks share a common architectural gap: they produce sub-tasks and sub-agents, but do not preserve the decomposition's causal history. When a sub-agent fails or produces a contested output, the ability to trace that sub-task back to the parent agent that decomposed it — and through that parent to the originating principal — is essential for both accountability and error diagnosis. The Recursive Spawning model's immutable parent chain directly addresses this gap, making hierarchical decomposition not only a planning mechanism but a governance structure.

\subsection{The BX3 Framework as Substrate for Immutable Parent Chains}

The BX3 Framework \cite{bx3framework} defines a universal architecture of functional roles within multi-agent systems, where every agent operates within a defined role context (Executor, Strategist, Critic, Coordinator, or Integrator) and communicates through typed role interfaces. Unlike monolithic agent architectures, BX3 treats role as a first-class constraint: an agent's capabilities, access scope, and delegation authority are determined by its role assignment within the hierarchy.

The BX3 Framework provides three structural properties that make immutable parent chains feasible. First, \textbf{role-bounded delegation}: an agent may only spawn sub-agents within its authorized scope; a Critic cannot spawn a Strategist, and an Executor cannot spawn outside its assigned task context. This prevents unbounded chain growth and ensures that every spawning event is role-consistent. Second, \textbf{immutable ancestry}: BX3 enforces that once an agent is spawned, its parent reference is cryptographically sealed and never re-parented. The parent chain forms a DAG from the root principal down to the leaf executor, and that DAG is append-only. Third, \textbf{forensic linkage}: every agent's execution record includes a parent chain hash that can be verified independently, enabling post-hoc audit of the full decomposition chain without relying on the agent system's internal state.

These properties address the accountability gaps identified in both mutable delegation systems and flat decomposition models. When combined with the provenance standards described in VAP, WCA, and TIBET, immutable parent chains provide a structural complement to evidentiary telemetry — ensuring that provenance records reference a chain that cannot be retrospectively altered. The Recursive Spawning model, introduced in the following section, builds on this substrate to define the precise mechanics of spawning, inheritance, and termination within the BX3 architecture.

% Bibliography
% Placeholder for compiled .bbl — populate via arxiv submission or manual BibTeX
% The following .bib keys match the citations above and should be included in the main .bib file:
% agentlifecycle, agentrm, altk, mcp, a2a, agenthive, provegent, forbesprov, vap, wca, tibet,
% gaat, fixeddelegation, mutabledelegation, stackplanner, ht n, adapt, reactree, hiplan,
% structuredagent, tdp, himac, bx3framework

\begin{thebibliography}{99}
% Entries will be populated from the .bib file during compilation
\end{thebibliography}